% File: Volume-II-Applications/appendices/appF-cosmological-perturbations.tex

\chapter{Cosmological Perturbation Theory}
\label{app:cosmo-perturb}

This appendix provides the technical background for linear cosmological
perturbation theory in the context of the lamphrodynamic fields of
\rsvpabbr{}.  Our primary objective is to formulate perturbations of the
triple $(\PhiField,\vField,\SField)$ on a homogeneous and isotropic
background, to identify gauge--invariant combinations, and to derive the
linearized evolution equations used in Chapters~4--7 of this volume.
The presentation parallels standard FLRW cosmology (Weinberg, Dodelson),
but the dynamical variables differ substantially.

\section{Background FLRW Geometry}

Let $(\mathcal M,g)$ be a spatially homogeneous and isotropic spacetime
with metric
\begin{equation}
\label{eq:flrw-metric}
\dd s^2 = -\dd t^2 + a(t)^2\gamma_{ij}\dd x^i\dd x^j,
\end{equation}
where $a(t)$ is the scale factor and $\gamma_{ij}$ is a maximally
symmetric spatial metric of constant curvature $k\in\{-1,0,+1\}$.  In
\rsvpabbr{}, the lamphrodynamic fields
\[
(\PhiField(t),\;\vField(t),\;\SField(t))
\]
satisfy the background lamphrodynamic equations derived in
\cref{chap:friedmann-lamph}.  Throughout this appendix, an overbar
denotes background quantities, e.g.\ $\overline\PhiField(t)$.

\section{Perturbation Ansatz}

We write the perturbed fields as
\begin{equation}
\begin{aligned}
\PhiField(t,\mathbf x) &= \overline\PhiField(t) + \delta\Phi(t,\mathbf x),\\
\vField_i(t,\mathbf x) &= \overline\vField_i(t) + \delta v_i(t,\mathbf x),\\
\SField(t,\mathbf x)  &= \overline\SField(t)  + \delta S(t,\mathbf x),
\end{aligned}
\end{equation}
where $\delta\Phi$, $\delta v_i$, and $\delta S$ are small fluctuations.

For the metric, we adopt the standard scalar--vector--tensor (SVT)
decomposition.  Let
\begin{equation}
\label{eq:metric-perturb}
\dd s^2 = -(1+2A)\dd t^2 + 2a(t)\bigl(B_{,i}+B_i\bigr)\dd t\,\dd x^i
+ a(t)^2\bigl[(1+2H_L)\gamma_{ij}+2H_{T\,ij}\bigr]\dd x^i\dd x^j,
\end{equation}
where scalar, vector, and tensor parts satisfy the usual transverse and
traceless conditions:
\[
B^i{}_{,i}=0,\quad H_{T\,ij}{}^{,j}=0,\quad H_{T\,i}{}^i = 0.
\]

\section{Gauge Transformations}

An infinitesimal coordinate transformation
$x^\mu\mapsto x^\mu+\xi^\mu$ induces transformations of both the metric
perturbations and the lamphrodynamic fields.  Decompose
$\xi^\mu=(\xi^0,\xi^{,i}+\xi^i_\perp)$ with $\xi^i_\perp{}_{,i}=0$.

\begin{lemma}[Gauge Transformations of Lamphrodynamic Perturbations]
\label{lem:gauge-lamphro}
Under an infinitesimal coordinate transformation,
\[
\delta\Phi \mapsto \delta\Phi - \dot{\overline\Phi}\,\xi^0,\qquad
\delta S \mapsto \delta S - \dot{\overline S}\,\xi^0,\qquad
\delta v_i \mapsto \delta v_i - \dot{\overline v}_i\,\xi^0,
\]
where overdot denotes derivative with respect to cosmological time.
\end{lemma}

\begin{proof}
Direct computation from the transformation law of scalar and vector
fields under infinitesimal diffeomorphisms, using background
homogeneity.
\end{proof}

Gauge--invariant combinations follow by combining these with metric
perturbations in analogy with the Bardeen potentials.  For instance,
define
\[
\PhiGI := \delta\Phi - \dot{\overline\Phi}\,\frac{1}{a}
\int^t a A\,\dd t',
\]
with analogous definitions for $\SGI$ and $\vGI_i$.  Tensor modes
$H_{T\,ij}$ are already gauge--invariant.

\section{Linearized Lamphrodynamic Equations}

The lamphrodynamic field equations (Volume~I, Chapters~5--8) can be
expanded to first order about the FLRW background.  Denote $\mathcal
E[\PhiField,\vField,\SField,g]=0$ schematically.  Linearizing yields
\[
\mathcal L_{\rm lin}
\begin{pmatrix}
\delta\Phi \\
\delta v_i\\
\delta S\\
A,B,H_L,H_{T\,ij}
\end{pmatrix}
=0,
\]
where $\mathcal L_{\rm lin}$ is a coupled, first--order (in time)
differential operator.  For scalar modes, a representative pair of
equations takes the form
\begin{align}
\delta\dot\Phi + \alpha_1 H\,\delta\Phi + \alpha_2\,\delta S
+ \alpha_3 A &= 0, \\
\delta\dot S + \beta_1 H\,\delta S + \beta_2\,\delta\Phi
+ \beta_3\nabla^2\delta\Phi &= 0,
\end{align}
where $H=\dot a/a$ is the Hubble parameter and $\alpha_i,\beta_i$ are
coefficients determined by the background solution and the coupling
structure in the lamphrodynamic action.

The full linearized system includes similar equations for $\delta v_i$
and couples to the metric perturbations via the modified Einstein
equations of Chapter~1.

\section{Scalar, Vector, and Tensor Modes}

Decompose perturbations of $\delta v_i$ as
\begin{equation}
\delta v_i = \delta v_{,i} + \delta v_i^{T},\qquad
\delta v_i^{T}{}^{,i}=0.
\end{equation}
Scalar modes couple to $(\delta\Phi,\delta S)$, vector modes evolve
independently at linear order, and tensor modes obey the modified
gravitational wave equation
\begin{equation}
\label{eq:tensor-eom}
\ddot H_{T\,ij} + 3H\,\dot H_{T\,ij} - \frac{\nabla^2}{a^2}H_{T\,ij}
= \Theta_{ij},
\end{equation}
where $\Theta_{ij}$ encodes entropy--induced anisotropic stress
(see Chapter~1).

\section{Power Spectra and Transfer Functions}

Define Fourier transforms by
\[
\delta\Phi(t,\mathbf k)=\int\dd^3x\,e^{-i\mathbf k\cdot\mathbf x}\,
\delta\Phi(t,\mathbf x),
\]
and similarly for $\delta S$ and $\delta v_i$.  The auto-- and
cross--power spectra are defined by
\[
\langle \delta\Phi(\mathbf k)\,\delta\Phi(\mathbf k')\rangle
= (2\pi)^3\delta(\mathbf k+\mathbf k')P_{\Phi\Phi}(k),
\]
with analogous definitions for $P_{SS}$, $P_{\Phi S}$, and the
velocity cross spectra.

The transfer functions $T_\Phi(k,t)$, $T_S(k,t)$, and $T_v(k,t)$ are
defined by
\[
\delta\Phi(k,t)=T_\Phi(k,t)\,\delta\Phi(k,t_i),
\]
etc., where $t_i$ denotes an early initial time.  These functions enter
predictions for structure formation (Chapter~6) and CMB anisotropies
(Chapter~8).

\section{Initial Conditions}

Initial conditions may be specified via quantum fluctuations in the
pre--lamphrodynamic epoch (Volume~III, Chapter~10) or by phenomenological
ans\"atze analogous to adiabatic and isocurvature modes.  For the
adiabatic case,
\[
\frac{\delta S}{\dot{\overline S}}
= \frac{\delta\Phi}{\dot{\overline\Phi}}
= \frac{\delta v_i}{\dot{\overline v}_i},
\]
generalizing the standard condition that all components share the same
initial phase.

\section{Remarks and Further Reading}

Standard introductions to cosmological perturbation theory include
Weinberg (\emph{Cosmology}) and Dodelson (\emph{Modern Cosmology}).
For the tensor decomposition and gauge--invariant formalism, see also
Kodama--Sasaki and Mukhanov.  The adaptation to \rsvpabbr{} follows by
linearizing the lamphrodynamic field equations and replacing the
standard stress--energy sources with the entropy--vector couplings of
Chapters~1--3.

